% Segunda Lei de Kepler.
% 25/05/2004 at 16:00h.
% Gambelas, Faro
% Written by Tordar

\documentclass[12pt,a4paper]{article}

\pagestyle{empty} 

\usepackage{graphicx}
\usepackage{marvosym}
\usepackage{pst-all} 
\usepackage{pst-slpe}

%Let's go: 

\begin{document} 
\null\vskip4cm
\begin{pspicture}(-2,-2)(2,2)
\psset{unit=10mm}
\savedata{\orbita}[{ 
{    10.0000,         0},
{     9.8582,    1.0864},
{     9.4488,    2.1431},
{     8.7639,    3.1368},
{     7.7911,    4.0263},
{     6.5122,    4.7558},
{     4.9046,    5.2409},
{     2.9484,    5.3421},
{     0.6617,    4.8005},
{    -1.7058,    3.1148},
{    -2.9006,   -0.1827},
{    -1.5197,   -3.3539},
{     0.8467,   -4.9180},
{     3.0942,   -5.4062},
{     5.0140,   -5.2839},
{     6.5924,   -4.7916},
{     7.8486,   -4.0611},
{     8.8033,   -3.1729},
{     9.4734,   -2.1810},
{     9.8700,   -1.1258},
{    10.0000,   -0.0399},}]
\dataplot[plotstyle=curve,linewidth=2pt,linecolor=black,linestyle=dashed,dash=4mm 1mm]{\orbita}
\pspolygon[linewidth=0mm,linecolor=lightgray,fillstyle=solid,fillcolor=lightgray](0,0)(9.9614,-0.6389)(10,0)(9.9614,0.6389)
\pscurve[linewidth=1mm,linecolor=blue]{->}(9.9614,-0.6389)(10,0)(9.9614,0.6389)
\pspolygon[linewidth=0mm,linecolor=lightgray,fillstyle=solid,fillcolor=lightgray](0,0)
(-1.8574,2.9300)
(-1.9959,2.7447)
(-2.1281,2.5507)
(-2.2530,    2.3480)
(-2.3697,    2.1368)
(-2.4773,    1.9173)
(-2.5749,    1.6900)
(-2.6614,    1.4556)
(-2.7360,    1.2146)
(-2.7979,    0.9679)
(-2.8463,    0.7166)
(-2.8806,    0.4617)
(-2.9004,    0.2045)
(-2.9054,   -0.0537)
(-2.8957,   -0.3117)
(-2.8713,   -0.5681)
(-2.8326,   -0.8216)
(-2.7801,   -1.0712)
(-2.7144,   -1.3156)
(-2.6364,   -1.5541)
(-2.5469,   -1.7858)
(-2.4468,   -2.0101)
(-2.3371,   -2.2265)
(-2.2187,   -2.4346)
(-2.0925,   -2.6342)
(-1.9595,   -2.8252)
(-1.8204,   -3.0074)
\pscurve[linewidth=1mm,linecolor=blue]{->}(-1.8574,2.9300)(-1.9959,2.7447)(-2.1281,2.5507)
(-2.2530,    2.3480)
(-2.3697,    2.1368)
(-2.4773,    1.9173)
(-2.5749,    1.6900)
(-2.6614,    1.4556)
(-2.7360,    1.2146)
(-2.7979,    0.9679)
(-2.8463,    0.7166)
(-2.8806,    0.4617)
(-2.9004,    0.2045)
(-2.9054,   -0.0537)
(-2.8957,   -0.3117)
(-2.8713,   -0.5681)
(-2.8326,   -0.8216)
(-2.7801,   -1.0712)
(-2.7144,   -1.3156)
(-2.6364,   -1.5541)
(-2.5469,   -1.7858)
(-2.4468,   -2.0101)
(-2.3371,   -2.2265)
(-2.2187,   -2.4346)
(-2.0925,   -2.6342)
(-1.9595,   -2.8252)
(-1.8204,   -3.0074)
\pscircle[linewidth=0mm,linecolor=white,fillstyle=solid,fillcolor=white]( 9.8582, 1.0864){0.45}
\pscircle[linewidth=0mm,linecolor=white,fillstyle=solid,fillcolor=white]( 9.8700,-1.1258){0.45}
\pscircle[linewidth=0mm,linecolor=white,fillstyle=solid,fillcolor=white](-1.7058, 3.1148){0.45}
\pscircle[linewidth=0mm,linecolor=white,fillstyle=solid,fillcolor=white](-1.5197,-3.3539){0.45}
\pscircle[linewidth=0mm,linecolor=white,fillstyle=solid,fillcolor=white](0,0){0.45}
\rput{0}( 9.8582, 1.0864){\scalebox{2.5}{\Mundus}}
\rput{0}( 9.8700,-1.1258){\scalebox{2.5}{\Mundus}} 
\rput{0}(-1.7058, 3.1148){\scalebox{2.5}{\Mundus}}
\rput{0}(-1.5197,-3.3539){\scalebox{2.5}{\Mundus}}
%\rput{0}(0,0){\scalebox{3.5}{$\odot$}}
\pscircle[linewidth=0.1mm,
linecolor=yellow,
fillstyle=ccslope,
slopebegin=red,
slopeend=yellow,
slopecenter=0.65 0.65,
slopesteps=25](0,0){0.7}  
\end{pspicture}
\end{document} 
